\documentclass[a4paper,10pt]{IEEEtran}

\title{Digital Circuits and VHDL -- Overview and Examples}
\author{Florian Hagemann\thanks{Thanks to \dots}}
\date{\today}




%======== Begin Document ==============

\begin{document}
\maketitle

\begin{abstract}
Here be an abstract\dots
\end{abstract}

\tableofcontents


%=========== Content ===================

\section{Introduction}

\section{Digital Circuits}
\IEEEPARstart{W}{hat} is a \emph{digital circuit}? To understand this, we need to examine each word separetely. A \emph{circuit} is a closed loop that allows electricity to flow. Typically it consists of a power source and some device (e.g. a motor, LED, etc.) that uses that electricity. A circuit needs to be closed in order for the electricity to flow through it \cite{wiktionary_circuit}. \emph{Digital} comes from the the Latin \textit{"digitalis"} meaning finger or toe. Digital can mean "something to do with digits (fingers or toes)" but can also mean the "property of representing values as discrete, often binary, numbers rather than a continuous spectrum" \cite{wiktionary_digital}.

A \textbf{digital circuit}

\subsection{Electronical Switches: Transistors}
\subsection{Making Computers Think}
\subsection{Examples}

\section{Introduction to VHDL}
\subsection{Hardware Description Languages}
\subsection{Basic Syntax}
\subsection{Blackbox Diagrams and Entities}
%etc...

\section{Combinatorical Circuits}
\subsection{}

\section{Sequential Circuits}

\section{Finite State Machines (FSM)}
\subsection{Moore Machines}
\subsection{Mealy Machines}


\nocite{*}
\bibliographystyle{IEEEtran}
\bibliography{sources}

\end{document}